\section{Discussion}

The results show that the evaluated dataset remains highly separable even after controlling visually equivalent samples through perceptual-hash grouping. The baseline MobileNetV2 model already achieved a high macro-F1 score of $0.9969 \pm 0.0052$, which indicates a performance saturation scenario. Under this condition, the margin for measurable improvement is small, and additional components do not necessarily produce statistically significant gains.

The strongest configuration was E2, which fused MobileNetV2 embeddings with explicit GLCM+DWT descriptors extracted from the original images. This model obtained the highest mean macro-F1 ($0.9993 \pm 0.0015$) and the lowest number of aggregate errors across five seeds. This result is consistent with previous works showing that handcrafted texture, color, and frequency descriptors can capture useful visual information for rice disease classification \cite{sutikno2026glcm_lbp_color,prity2026ffnn}. In this study, however, these descriptors were not evaluated as a standalone classifier input, but as a complementary branch fused with a lightweight CNN.

The segmentation-only configuration, E1, achieved the lowest inference time while maintaining high classification performance. This suggests that removing background regions can simplify the visual input and improve efficiency without substantially reducing discriminative information. This finding is aligned with segmentation-aware rice disease classification approaches, where region-focused processing helps the classifier attend to relevant leaf areas \cite{shakeel2026contour}. Nevertheless, in the present ablation, segmentation alone produced only a small average improvement over the baseline and did not reach statistical significance.

The proposed hybrid model, E3, was competitive but did not outperform E2 or show a statistically significant improvement over E0. This can be explained by three main factors. First, the baseline was already very strong, leaving little room for additional gains. Second, segmentation may remove not only irrelevant background but also some contextual or boundary information useful for classification. Third, the segmented CNN input and the texture descriptors may provide partially redundant information, limiting the benefit of their fusion. Therefore, E3 improves over the baseline in aggregate accuracy, but its added complexity does not translate into a clear performance advantage.

Compared with recent CNN-based rice disease classification studies, the results confirm that transfer learning backbones can achieve strong performance in this task \cite{ismail2026densenet}. However, direct numerical comparison with prior works should be interpreted cautiously because datasets, splits, preprocessing strategies, leakage controls, and evaluation protocols differ across studies. The main contribution of this paper is therefore not a claim of absolute superiority over previous models, but a controlled comparison of segmentation, explicit texture fusion, and their combination under the same experimental conditions.

Several threats to validity remain. First, the pHash-based split reduced the risk of exact or near-exact visual leakage by keeping images with the same perceptual hash within the same split. However, this procedure does not guarantee the removal of all visually or semantically similar samples, such as images captured under the same background, lighting, leaf position, or acquisition session. Second, the evaluation used one fixed pHash-controlled train/validation/test split and five random seeds. Therefore, the multi-seed evaluation measures training stability under the same data partition, while sensitivity to alternative dataset partitions was not evaluated. Third, no external field-acquired dataset was included; therefore, the reported performance should be interpreted within the conditions of the evaluated public dataset and not as direct evidence of generalization to uncontrolled agricultural environments. Finally, the segmentation method was based on classical L*a*b* Otsu thresholding, which is simple and reproducible, but may be less robust when images contain complex backgrounds, shadows, overlapping leaves, or lighting variations.

Overall, the evidence suggests that explicit texture fusion was the most effective enhancement in this experimental setting, segmentation was the most efficient configuration, and the proposed segmentation-texture hybrid was competitive but not statistically superior. Future work should validate the same ablation design on external field-acquired datasets to determine whether segmentation and texture fusion provide larger benefits under stronger domain shift.
